\documentclass[8pt, a4paper, landscape]{extarticle}

% Компактные поля
\usepackage[margin=1cm]{geometry}
\usepackage[T2A]{fontenc}
\usepackage[utf8]{inputenc}
\usepackage[english, russian]{babel}
\usepackage{amsmath, amssymb, amsfonts}
\usepackage{enumitem}
\usepackage{multicol} % Для колонок
\usepackage[most]{tcolorbox}

% Настройка списка для экономии места
\setlist[itemize]{leftmargin=*, noitemsep, topsep=0pt}
\setlist[enumerate]{leftmargin=*, noitemsep, topsep=0pt}

% Стиль блока билета
\newtcolorbox{ticket}[1]{
    colback=white,
    colframe=black!70,
    coltitle=white,
    fonttitle=\bfseries\small,
    title={#1},
    arc=0mm,
    boxrule=0.5pt,
    left=1mm, right=1mm, top=0mm, bottom=0mm,
    after skip=2mm,
    middle=0.3mm % <--- Уменьшает отступ вокруг разделительной линии
}

\newcommand{\Ra}{\Rightarrow}
\newcommand{\N}{\mathbb{N}}
\newcommand{\Q}{\mathbb{Q}}
\newcommand{\Z}{\mathbb{Z}}
\newcommand{\R}{\mathbb{R}}
\newcommand{\eps}{\varepsilon}

\DeclareMathOperator{\sign}{sign}

\begin{document}

% Заголовок не нужен для шпоры, сразу к делу
\begin{multicols*}{3} % 3 колонки для компактности

% ================= РАЗДЕЛ 1 =================

\begin{ticket}{1.1 Комплексные числа. Эйлер}
\textbf{Опр:} $z = x+iy$, $i^2=-1$. Сравнивать ($>, <$) нельзя.
\textbf{Формы:} Алгебр ($x+iy$), Тригон ($r(\cos\phi + i\sin\phi)$).
\textbf{Эйлер:} $e^{ix} = \cos x + i\sin x$.
\textbf{Опр $e^z$:} $e^{x+iy} = e^x(\cos y + i\sin y)$.
\textit{Док-во Эйлера:} Разложить $e^{ix}$ в ряд Тейлора $\to$ сгруппировать слагаемые с $i$ и без $\to$ получить ряды $\cos x$ и $\sin x$.
\end{ticket}

\begin{ticket}{1.2 Муавр. Корни}
\textbf{Муавр:} $(\cos \phi + i\sin \phi)^n = (\cos n\phi + i\sin n\phi)$.
\textbf{Корни:} $\sqrt[n]{z}$ имеет $n$ значений. Точки — вершины правильного $n$-угольника.
\textbf{Формула:} $\omega_k = \sqrt[n]{r} \cdot e^{i \frac{\phi + 2\pi k}{n}}$.
\end{ticket}

\begin{ticket}{1.3 Неравенство треугольника}
\textbf{Теорема:} $|z_1 + z_2| \le |z_1| + |z_2|$.
\textbf{Док-во:}
1. Возвести $|z_1+z_2|^2$ в квадрат через сопряженные: $(z_1+z_2)(\overline{z_1}+\overline{z_2})$.
2. Раскрыть: $|z_1|^2 + |z_2|^2 + 2\text{Re}(z_1\overline{z_2})$.
3. Оценка: $\text{Re}(w) \le |w| \Ra \le (|z_1|+|z_2|)^2$.
\textbf{Следствие:} $||z_1| - |z_2|| \le |z_1 - z_2|$.
\textit{Док-во:} Записать $z_1 = (z_1-z_2) + z_2$, применить прямую теорему.
\end{ticket}

\begin{ticket}{1.4 Свойства операций (Множества)}
\textbf{Дистрибутивность:} $A \cap (B \cup C) = (A \cap B) \cup (A \cap C)$.\\
$A + (BC) = (A + B)(A + C)$
\textbf{Док-во:} Через включение множеств $\subset$ и $\supset$. Перебор случаев принадлежности элемента $x$.
\end{ticket}

\begin{ticket}{1.5 ММИ. Бином Ньютона}
\textbf{ММИ:} 1. База ($n=1$). 2. Шаг ($n=k \to n=k+1$).
\textbf{Бином:} $(a+b)^n = \sum C_n^k a^{n-k}b^k$.
\textbf{Док-во:} Индукция.
Шаг: $(a+b)^{k+1} = (a+b)(a+b)^k$.
Ключ: $C_k^i + C_k^{i-1} = C_{k+1}^i$ (Треугольник Паскаля).
\end{ticket}

\begin{ticket}{1.6 Средние. Обратная индукция}
\textbf{Теорема (Коши):} $\frac{\sum a_i}{n} \ge \sqrt[n]{\prod a_i}$.
\textbf{Док-во (Обратная индукция):}
1. \textbf{Вверх:} Доказать для $n=2$ (Через $(a-b)^2 \ge 0$), потом $4, 8 \dots 2^k$ через предположение, что для k верно и доказанное свойство для n=2.
2. \textbf{Вниз:} Если верно для $n$, верно для $n-1$.
Ключ: Взять $a_n = \text{ср арифм первых } n-1$. Привести, что дробь равна $a_n$ и больше ср. геом. возвести в степень, поделить, взять корень, поделить. 
\end{ticket}

% ================= РАЗДЕЛ 2: ДЕЙСТВИТЕЛЬНЫЕ ЧИСЛА =================

% ================= БИЛЕТ 2.1 =================
\begin{ticket}{2.1 Дедекиндовы сечения. Опр. числа}
\textbf{Проблема $\Q$:} $\sqrt{2} \notin \Q$. Есть «дырки».
\textbf{Сечение:} Последовательность вложенных отрезков с рац. концами $I_n = [a_n, b_n]$, длина которых $\to 0$.
\textbf{Классы:}
1. Левый: $r <$ хотя бы одного $a_n$.
2. Правый: $r >$ хотя бы одного $b_n$.
3. Центр: $\{r \in \Q \mid \forall n: a_n \le r \le b_n\}$ (0 или 1 число).
\textbf{Опр. $\R$:} Действительное число — это класс эквивалентных последовательностей отрезков.
\tcblower
\textbf{Критерий эквивалентности (Док-во):}
$\{I_n\} \sim \{I'_n\} \iff \lim(a_n - a'_n) = 0$.
\textbf{($\Ra$)}: От противного. Пусть $\lim \ne 0$. Тогда $\exists \eps$: отрезки $I_n$ и $I'_n$ разойдутся на расстояние $\eps$.
Т.к. длины $\to 0$, найдется место между ними.
Вставим туда число $r$. Оно попадет в ПРАВЫЙ класс для одних и в ЛЕВЫЙ для других $\Ra$ классы разные. Противоречие.
\textbf{($\Leftarrow$)}: Пусть $\lim = 0$. Если классы разные, то найдется $r$ между отрезками. Но тогда расстояние между ними $>0$, а должно стремиться к 0.
\end{ticket}

% ================= БИЛЕТ 2.2 =================
\begin{ticket}{2.2 Лемма об отделимости}
\textbf{Условие:} $A, B \subset \R$, оба не пустые.
$\forall a \in A, \forall b \in B \implies a \le b$.
\textbf{Утв:} $\exists c \in \R: \forall a \in A, \forall b \in B \implies a \le c \le b$.
\tcblower
\textbf{Доказательство:}
\textbf{Случай 1:} $A \cap B \ne \varnothing$.
Пусть $c \in A \cap B$. Тогда $c \in A \Ra c \le b$. И $c \in B \Ra a \le c$. Всё доказано.
\textbf{Случай 2:} $A \cap B = \varnothing$ (значит $a < b$ строго).
1. Расширим множества до покрытия всей прямой:
$B' = \{y \mid \exists b \in B: b \le y\}$ (всё, что правее $B$).
$A' = \R \setminus B'$. Тогда $\forall x \in A', \forall y \in B' \Ra x < y$.
2. Строим разделяющую точку (метод деления пополам):
Берем $a_1 \in A', b_1 \in B'$. Делим пополам $m = \frac{a_1+b_1}{2}$.
Если $m \in A'$, то $a_2=m$. Если $m \in B'$, то $b_2=m$.
3. Получаем вложенные стягивающиеся отрезки $[a_n, b_n]$.
4. По аксиоме/свойству $\R$, у них есть общая точка $c$.
Она и будет разделять множества.
\end{ticket}

% ================= БИЛЕТ 2.3 =================
\begin{ticket}{2.3 Теорема Вейерштрасса (Sup/Inf)}
\textbf{Опр:} $M = \sup E$ (точная верхняя грань), если:
1. $\forall x \in E: x \le M$ (верхняя грань).
2. $\forall M' < M \exists x \in E: x > M'$ (наименьшая).
\textbf{Теорема:} Если $E \ne \varnothing$ и ограничено сверху, то $\exists \sup E$.
\tcblower
\textbf{Доказательство (через отделимость):}
1. Пусть множество $A = E$.
2. Пусть множество $B = \{ \text{все верхние грани множества } E \}$.
Т.к. $E$ ограничено сверху, $B \ne \varnothing$.
3. Проверим условие леммы: берем любое $a \in E$ и любой $b \in B$ (верхнюю грань). По определению грани $a \le b$.
4. По Лемме об отделимости $\exists c: a \le c \le b$.
5. Докажем, что $c = \sup E$:
- Т.к. $a \le c$ для всех $a \in E$, то $c$ — верхняя грань ($c \in B$).
- Т.к. $c \le b$ для всех $b \in B$, то $c$ — самая маленькая из верхних граней.
\end{ticket}

% ================= БИЛЕТ 2.4 =================
\begin{ticket}{2.4 Принцип Кантора (Вложенные отрезки)}
\textbf{Теорема:} Пусть $[a_1, b_1] \supset [a_2, b_2] \supset \dots$
Тогда пересечение $\bigcap [a_k, b_k] \ne \varnothing$.
Если длина $\to 0$, то точка единственная.
\tcblower
\textbf{Доказательство (через Вейерштрасса):}
1. Рассмотрим множество левых концов $A = \{a_1, a_2, \dots\}$.
2. Оно ограничено сверху (любым правым концом, например $b_1$, т.к. $a_n \le b_n \le b_1$).
3. По Т. Вейерштрасса $\exists c = \sup \{a_n\}$.
4. Так как $c$ — супремум левых концов, то $a_n \le c$.
5. Докажем, что $c \le b_n$.
Предположим $c > b_k$. Тогда $b_k$ не верхняя грань для $A$.
Значит $\exists a_m > b_k$.
Возьмем $N = \max(m, k)$. Т.к. отрезки вложены:
$a_N \ge a_m > b_k \ge b_N \implies a_N > b_N$. Противоречие.
Итог: $a_n \le c \le b_n$, точка $c$ лежит во всех отрезках.
\end{ticket}

% ================= БИЛЕТ 2.5 =================
\begin{ticket}{2.5 Полнота по Дедекинду (из Кантора)}
\textbf{Теорема:} Всякое сечение $(L, R)$ в $\R$ определяется числом $c$ (границей), таким что $L \le c \le R$. (В $\R$ нет дырок).
\tcblower
\textbf{Доказательство:}
1. Возьмем $x_0 \in L$ и $y_0 \in R$.
2. Строим отрезки делением пополам. $m = \frac{x_0+y_0}{2}$.
- Если $m \in L$, то $x_1 = m, y_1 = y_0$.
- Если $m \in R$, то $x_1 = x_0, y_1 = m$.
3. Получили вложенные отрезки $[x_n, y_n]$, где левый конец всегда в $L$, правый в $R$. Длина $\to 0$.
4. По Принципу Кантора $\exists ! c$, принадлежащее всем отрезкам.
5. Это $c$ и есть граница.
Любое $l \in L$ меньше $c$ (иначе попадет в правые концы).
Любое $r \in R$ больше $c$.
\end{ticket}

% ================= БИЛЕТ 2.6 =================
\begin{ticket}{2.6 Несчетность $\R$}
\textbf{Счетное:} Можно занумеровать (биекция с $\N$). $\Q$ — счетно.
\textbf{Теорема:} Отрезок $[0, 1]$ (и $\R$) — несчетен.
\tcblower
\textbf{Доказательство (Метод вложенных отрезков):}
От противного. Пусть все числа из $[0, 1]$ можно записать в список: $x_1, x_2, x_3 \dots$
1. Делим отрезок $[0, 1]$ на 3 части.
Выбираем ту часть $\Delta_1$, в которую НЕ попала точка $x_1$.
2. Делим $\Delta_1$ на 3 части.
Выбираем часть $\Delta_2$, в которую НЕ попала точка $x_2$.
3. Продолжаем. Получаем вложенные отрезки $\Delta_k$.
В $\Delta_k$ нет точки $x_k$.
4. По Т. Кантора $\exists c \in \bigcap \Delta_k$.
5. Это число $c$ существует, но оно не равно $x_1$ (т.к. $c \in \Delta_1$), не равно $x_2$ (т.к. $c \in \Delta_2$) и т.д.
6. Мы нашли число, которого нет в нашем "полном" списке. Противоречие.
\end{ticket}



% ================= БИЛЕТ 3.1 =================
\begin{ticket}{3.1 Свойства сходящихся}
\textbf{Опр:} $\lim a_n = a \iff \forall \eps>0 \exists N: \forall n>N |a_n-a|<\eps$.
\textbf{Т1:} Постоянная $a_n=a \Ra \lim a_n=a$.
\textbf{Т2 (Единственность):} Предел единственен.
\textbf{Т3 (Ограниченность):} Сходящаяся $\Ra$ Ограниченная.
\tcblower
\textbf{Док-во единственности:}
От противного. Пусть два предела $A \ne B$.
Взять $\eps = |A-B|/2$. Окрестности не пересекаются.
Начиная с $N = \max(N_1, N_2)$ точки должны быть в обеих. Противоречие.
\textbf{Док-во ограниченности:}
Взять $\eps=1$. С $N$ все точки в $(a-1, a+1)$.
Вне окрестности только конечное число точек $a_1 \dots a_N$.
$M = \max(|a_1| \dots |a_N|, |a|+1)$.
\end{ticket}

% ================= БИЛЕТ 3.2 =================
\begin{ticket}{3.2 Предельный переход}
\textbf{Теорема:} Если $a_n \to a, b_n \to b$ и $a_n \le b_n$ (начиная с некоторого $N$), то $a \le b$.
\textbf{Важно:} Строгое неравенство $a_n < b_n$ может перейти в нестрогое $a \le b$.
\tcblower
\textbf{Идея доказательства:}
От противного. Пусть $a > b$.
Взять $\eps = (a-b)/2$.
Окрестности $U(a)$ и $U(b)$ не пересекаются.
$U(b)$ левее $U(a)$.
$b_n$ попадет в левую, $a_n$ в правую $\Ra b_n < a_n$. Противоречие с условием.
\end{ticket}

% ================= БИЛЕТ 3.3 =================
\begin{ticket}{3.3 Зажатая последовательность}
\textbf{Теорема (о двух милиционерах):}
Если $a_n \le b_n \le c_n$ и $\lim a_n = \lim c_n = A$, то $\lim b_n = A$.
\tcblower
\textbf{Идея доказательства:}
По определению.
Для $\eps$ найти $N_1$ для $a_n$ и $N_2$ для $c_n$.
Взять $N = \max$.
$A-\eps < a_n \le b_n \le c_n < A+\eps$.
Значит $b_n \in (A-\eps, A+\eps)$.
\end{ticket}

% ================= БИЛЕТ 3.4 =================
\begin{ticket}{3.4 Сохранение знака / Модуль}
\textbf{Т1 (Знак):} Если $a_n \to a \ne 0$, то начиная с некоторого $N$, $\sign a_n = \sign a$.
\textbf{Т2 (Модуль):} $|a_n| \to |a|$.
\tcblower
\textbf{Док-во знака:}
Пусть $a > 0$. Взять $\eps = a/2$.
Тогда $a_n > a - a/2 = a/2 > 0$.
\textbf{Док-во модуля:}
Неравенство треугольника: $||a_n| - |a|| \le |a_n - a| < \eps$.
\end{ticket}

% ================= БИЛЕТ 3.5 =================

\begin{ticket}{3.5 Б.м. последовательности}
\textbf{Опр:} $\alpha_n \to 0$ (б.м.).
\textbf{Т (Связь):} $a_n \to a \iff a_n = a + \alpha_n$.
\textbf{Т (Свойства б.м.):}
1. Сумма/разность б.м. — б.м. ($\alpha_n \pm \beta_n \to 0$).
2. Произведение б.м. на \textbf{ограниченную} — б.м.
\tcblower
\textbf{Док-во Т:}
($\Ra$) $a_n \to a \Ra |a_n - a| < \eps$. Обозначим $\alpha_n = a_n - a$. Тогда $|\alpha_n| < \eps \Ra \alpha_n \to 0$.
($\Leftarrow$) $a_n = a + \alpha_n$. Тогда $|a_n - a| = |\alpha_n| < \eps \Ra a_n \to a$.

\textbf{Док-во (Сумма):}
Взять $\eps$. Для $\alpha_n$ найти $N_1$ (для $\eps/2$). Для $\beta_n$ найти $N_2$ (для $\eps/2$).
Для $n > \max(N_1, N_2)$:
$|\alpha_n \pm \beta_n| \le |\alpha_n| + |\beta_n| < \eps/2 + \eps/2 = \eps$.

\textbf{Док-во (Произведение):}
$\beta_n$ ограничена $\Ra \exists M: |\beta_n| \le M$.
Для $\alpha_n$ найти $N$ для $\eps/M$.
Тогда $|\alpha_n \beta_n| = |\alpha_n| \cdot |\beta_n| < (\eps/M) \cdot M = \eps$.
\end{ticket}

% ================= БИЛЕТ 3.6 =================
\begin{ticket}{3.6 Бесконечно большие (б.б.)}
\textbf{Опр:} $\forall E>0 \exists N \forall n>N: |a_n| > E$.
\textbf{Связь:}
1. $a_n \to \infty (a_n \ne 0) \Ra 1/a_n \to 0$.
2. $\alpha_n \to 0 (\alpha_n \ne 0) \Ra 1/\alpha_n \to \infty$.
\tcblower
\textbf{Идея доказательства:}
Меняем местами $\eps$ и $1/E$.
$|a_n| > E \iff 1/|a_n| < 1/E$.
Пусть $\eps = 1/E$.
\end{ticket}

% ================= БИЛЕТ 3.7 =================
\begin{ticket}{3.7 Арифметика пределов}
Если $a_n \to a, b_n \to b$:
1. $a_n \pm b_n \to a \pm b$.
2. $a_n b_n \to ab$.
3. $a_n/b_n \to a/b$ (при $b \ne 0$).
\tcblower
\textbf{Идея доказательства:}
Представить $a_n = a + \alpha_n, b_n = b + \beta_n$.
Для умножения:
$a_n b_n = (a+\alpha)(b+\beta)$.
В скобке — сумма б.м. (так как $a, b$ константы $\to$ ограничены).
Для деления: доказать, что $1/b_n$ ограничена (отделена от нуля).
Так как $b \ne 0$, возьмем $\varepsilon = \frac{|b|}{2}$.
Начиная с некоторого $N$: $|b_n - b| < \frac{|b|}{2}$.
По нер-ву треугольника: $|b| - |b_n| < \frac{|b|}{2} \Ra |b_n| > \frac{|b|}{2}$.
Тогда $\left|\frac{1}{b_n}\right| < \frac{2}{|b|}$.
Значит, дробь $\frac{1}{b b_n}$ ограничена константой.
Произведение б.м. на ограниченную есть б.м.
\end{ticket}

% ================= БИЛЕТ 3.8 =================
\begin{ticket}{3.8 Монотонные последовательности}
\textbf{Опр:} $\{a_n\} \uparrow$ (неубыв), если $a_n \le a_{n+1}$.
$\{a_n\} \downarrow$ (невозр), если $a_n \ge a_{n+1}$.
\textbf{Т (Вейерштрасс):}
1. $\{a_n\} \uparrow$ и огр. сверху $\Ra \lim a_n = \sup a_n$.
2. $\{a_n\} \downarrow$ и огр. снизу $\Ra \lim a_n = \inf a_n$.
\textbf{Критерий:} Монотонная сходится $\iff$ она ограничена.
\tcblower
\textbf{Док-во (для возрастающей):}
Пусть $A = \sup a_n$.
1. $\forall n: a_n \le A$ (верхняя грань).
2. $\forall \eps>0 \exists N: a_N > A-\eps$ (точная грань).
Т.к. $a_n \uparrow$, то $\forall n>N: a_n \ge a_N$.
Итог: $A-\eps < a_N \le a_n \le A < A+\eps$.
Значит $|a_n - A| < \eps$.
\textbf{Если неограничена:}
$\{a_n\} \uparrow$ и не огр. сверху $\Ra \forall E \exists N: a_N > E \Ra \forall n>N: a_n > E \Ra a_n \to +\infty$.
\end{ticket}

% ================= БИЛЕТ 3.9 =================
\begin{ticket}{3.9 Число e}
\textbf{Опр:} $e = \lim_{n\to\infty} (1+\frac{1}{n})^n$.
\textbf{Док-во существования:}
Нужно доказать, что $x_n = (1+1/n)^n$ возрастает и ограничена.
1. \textbf{Возрастание:} Неравенство Коши для $n+1$ чисел ($1, 1+\frac{1}{n}, \dots$).
$\sqrt[n+1]{1 \cdot (1+\frac{1}{n})^n} < \frac{1 + n(1+\frac{1}{n})}{n+1} = \frac{n+2}{n+1} = 1+\frac{1}{n+1}$.
Возводим в $n+1$: $(1+\frac{1}{n})^n < (1+\frac{1}{n+1})^{n+1}$.
2. \textbf{Ограниченность:} Рассм. $y_n = (1+\frac{1}{n})^{n+1}$. Она убывает.
$x_n < y_n < y_1 = 4$.
\textbf{Связь с рядом:} $e = \sum_{k=0}^\infty \frac{1}{k!}$.
\end{ticket}

% ================= БИЛЕТ 3.10 =================
\begin{ticket}{3.10 Теорема Больцано-Вейерштрасса}
\textbf{Теорема:} Из любой \textbf{ограниченной} последовательности можно выделить \textbf{сходящуюся} подпоследовательность.
\textbf{Следствие:} Из любой (даже неогр.) послед-ти можно выделить подпоследовательность, имеющую предел (конечный или $\infty$).
\tcblower
\textbf{Док-во (Метод деления пополам):}
1. $a_n \in [A, B]$. Делим отрезок пополам.
2. Выбираем половину, где бесконечно много членов. Называем $[a_1, b_1]$.
3. Повторяем: $[a_k, b_k]$.
4. Получаем вложенные отрезки, длина $\to 0$.
5. По Лемме Кантора есть общая точка $c$.
6. Выбираем $x_{n_k} \in [a_k, b_k]$. По теореме о зажатой посл-ти $x_{n_k} \to c$.
\end{ticket}

% ================= БИЛЕТ 3.11 =================
\begin{ticket}{3.11 Частичные пределы}
\textbf{Опр:} $a$ — частичный предел (ч.п.), если $\exists a_{n_k} \to a$.
\textbf{Критерий:} $a$ — ч.п. $\iff$ в любой $\eps$-окрестности $a$ бесконечно много членов $a_n$.
\tcblower
\textbf{Док-во ($\Ra$):}
$a_{n_k} \to a \Ra \forall \eps \exists K: \forall k>K a_{n_k} \in U_\eps(a)$.
Индексов $k$ бесконечно много $\Ra$ членов бесконечно много.
\textbf{Док-во ($\Leftarrow$):}
Строим $a_{n_k}$ индуктивно.
В $(a-1, a+1)$ берем $a_{n_1}$.
В $(a-1/2, a+1/2)$ берем $a_{n_2}$ с $n_2 > n_1$.
В $(a-1/k, a+1/k)$ берем $a_{n_k}$ с $n_k > n_{k-1}$.
\end{ticket}

% ================= БИЛЕТ 3.12 =================
\begin{ticket}{3.12 Критерий Коши (Последовательности)}
\textbf{Фунд-ть:} $\forall \eps \exists N \forall n, m > N: |a_n - a_m| < \eps$.
\textbf{Теорема:} Сходится $\iff$ Фундаментальна.
\tcblower
\textbf{($\Ra$):} Нер-во $\triangle$. $|a_n - a_m| \le |a_n - a| + |a - a_m| < \eps/2 + \eps/2$.
\textbf{($\Leftarrow$):}
1. Фунд. $\Ra$ Ограничена (взять $\eps=1$).
2. Ограничена $\Ra \exists$ сходящаяся $a_{n_k} \to A$ (Б-В).
3. Доказать, что вся $a_n \to A$.
$|a_n - A| \le |a_n - a_{n_k}| + |a_{n_k} - A|$.
Первое $<\eps/2$ (фунд.), второе $<\eps/2$ (предел).
\end{ticket}

% ================= БИЛЕТ 3.13 =================
\begin{ticket}{3.13 Верхний/Нижний пределы}
\textbf{Опр:} $\varlimsup a_n$ — наибольший ч.п., $\varliminf a_n$ — наименьший.
\textbf{Теорема:} У любой посл-ти существуют (конечные или $\infty$).
\tcblower
\textbf{Док-во существования верхнего предела ($L$):}
\textbf{Случай 1 (Неогр. сверху):}
$\forall k \exists a_{n_k} > k$. Эта подпоследовательность $\to +\infty$.
Значит $+\infty \in P$. $L = +\infty$.
\textbf{Случай 2 (Огр. сверху):}
Пусть $P^*$ — множество конечных ч.п.
\textbf{а) Если $P^* = \varnothing$:}
Докажем, что $a_n \to -\infty$.
Иначе $\exists E$ и подпоследовательность $a_{n_k} \ge -E$.
Она огр. (сверху условием случая 2, снизу числом $-E$).
По Б-В из неё можно выделить сходящуюся $\Ra$ есть конечный ч.п. Противоречие.
Раз $a_n \to -\infty$, то $L = -\infty$.
\textbf{б) Если $P^* \ne \varnothing$:}
$P^*$ ограничено сверху (т.к. $a_n$ ограничена).
По Вейерштрассу $\exists M = \sup P^*$.
Докажем, что $M \in P^*$ (т.е. супремум достижим).
В любой $\eps$-окрестности $M$ есть какой-то ч.п. $p \in P^*$.
В окрестности $p$ бесконечно много членов $a_n$.
Значит и в окрестности $M$ их бесконечно много.
Следовательно, $M$ — ч.п. и $L = M$.

\end{ticket}

% ================= БИЛЕТ 3.14 (Основы) =================
\begin{ticket}{3.14 Ряды. Коши. Сравнение}
\textbf{Опр:} $S_n = \sum_{k=1}^n a_k$. Ряд сходится $\iff \exists \lim S_n = S$.
\textbf{Кр. Коши:} $\sum a_n$ сход. $\iff$
$\forall \eps > 0 \exists N \forall n \ge N \forall p: |\sum_{k=n+1}^{n+p} a_k| < \eps$.
\textbf{Необх. условие:} $\sum a_k$ сход $\Ra a_n \to 0$ (взять $p=1$ в Коши).
\textbf{Абсолютная сходимость:} Сходится $\sum |a_n|$.
\textbf{Т:} Абсолютно $\Ra$ Условно (из нер-ва $\triangle$ в критерии Коши).
\tcblower
\textbf{Док-во Коши:}
Введем $X_n = \sum_{k=1}^n a_k$.
Тогда $\sum_{k=n+1}^{n+p} a_k = X_{n+p} - X_n$.
\textbf{($\Leftarrow$)}: Условие $\iff |X_{n+p} - X_n| < \eps$.
Значит $\{X_n\}$ — фундаментальна.
Фунд. $\Ra$ Сходится (Критерий Коши для посл.).
Сход. $\{X_n\}$ $\Ra$ Ряд сходится.
\textbf{($\Rightarrow$)}: Ряд сход. $\Ra \{X_n\}$ сход.
Сход. $\Ra$ Фунд. $\Ra |X_{n+p} - X_n| < \eps$.
\textbf{Признак сравнения ($0 \le a_n \le b_n$):}
1. $\sum b_n$ сход $\Ra \sum a_n$ сход.
2. $\sum a_n$ расход $\Ra \sum b_n$ расход.
\textit{Док-во:} $A_n \le B_n$.
Если $\sum b_n$ сход $\Ra B_n$ огр $\Ra A_n$ огр.
Т.к. $a_n \ge 0$, то $A_n \uparrow$.
Монотонная и огр. $\Ra$ сходится.
\end{ticket}

% ================= БИЛЕТ 3.15 (Признаки) =================
\begin{ticket}{3.15 Коши и Даламбер}
\textbf{Признак Коши:} $q = \varlimsup \sqrt[n]{|a_n|}$.
\textbf{Признак Даламбера:} $q = \lim |\frac{a_{n+1}}{a_n}|$.
\textbf{Вывод:}
$q < 1 \Ra$ Абсолютно сходится.
$q > 1 \Ra$ Расходится (нарушено необх. усл., $a_n \not\to 0$).
$q = 1 \Ra$ Неопределенность.
\tcblower
\textbf{Док-во (Коши):}
Сравнение с геом. прогрессией.
Если $q < 1$, берем $p \in (q, 1)$. Начиная с $N$: $\sqrt[n]{|a_n|} < p \Ra |a_n| < p^n$. Сходится ряд с геом прогрессией, которая больше ряда с $|a_n|$ значит ряд тоже сходится.
Док-во расходимости ($q > 1$):
Т.к. $\varlimsup > 1$, существует подпоследовательность $n_k$, такая что:
$\sqrt[n_k]{|a_{n_k}|} > 1 \implies |a_{n_k}| > 1$.
Если бесконечно много членов $>1$, то $\lim a_n \ne 0$.
Нарушено необх. условие $\Ra$ Расходится.
\textbf{Док-во (Даламбер):}
Если $q < 1$, берем $p \in (q, 1)$. $|a_{n+1}| < p|a_n|$.
По индукции $|a_n| < |a_N| \cdot p^{n-N} = C \cdot p^n$. Сходится.
Для $q>1$ $|a_{n+1}|>|a_n| \Rightarrow a_n \not\to 0$
\end{ticket}

% ================= БИЛЕТ 3.16(Конденсация) =================
\begin{ticket}{3.16 Конденсационный признак}
\textbf{Теорема:} Если $a_n \ge 0$ и $a_n \downarrow$, то
$\sum a_n \sim \sum 2^k a_{2^k}$ (сходимость равносильна).
\textbf{Исследование $\sum \frac{1}{n^p}$:}
\tcblower
\textbf{Идея доказательства:}
Оценка частичных сумм $S_n$ через $T_k = \sum 2^j a_{2^j}$.
1. Сверху: $a_2+a_3 \le 2a_2 \dots \Ra S_{2^{k+1}-1} \le a_1 + T_k$.
Если $T_k$ сход $\Ra S_n$ огр $\Ra$ сходц.
2. Снизу: $a_3+a_4 \ge 2a_4 \dots \Ra S_{2^{k+1}} \ge a_1 + a_2 + \frac{1}{2}T_{k+1}$.
Если $T_k \to \infty \Ra S_n \to \infty$.

\textbf{Исследование:}
Применим конд. признак. Сжатый ряд:
$\sum 2^k \cdot \frac{1}{(2^k)^p} = \sum 2^k \cdot 2^{-kp} = \sum (2^{1-p})^k$.
Это геом. прогрессия со знаменателем $Q = 2^{1-p}$.
\end{ticket}


% \begin{multicols*}{3}

% ================= БИЛЕТ 4.1-4.2 =================
\begin{ticket}{4.1-4.2 Определения предела функции}
\textbf{По Коши ($\varepsilon-\delta$):} $\lim_{x\to a} f(x)=A \iff$
$\forall \eps>0 \exists \delta>0 \forall x: 0 < |x-a| < \delta \implies |f(x)-A| < \eps$.
\textbf{По Гейне (посл-ти):} $\lim_{x\to a} f(x)=A \iff$
$\forall x_n \to a (x_n \ne a) \implies f(x_n) \to A$.
\tcblower
\textbf{Эквивалентность:}
\textbf{Коши $\Ra$ Гейне:} Составим $x_n \to a \Ra |x_n-a| < \delta \Ra |f(x_n)-A| < \eps$.
\textbf{Гейне $\Ra$ Коши:} (От противного).
Пусть Коши неверно: $\exists \eps \forall \delta \exists x_\delta \dots$
Берем $\delta = 1/n$. Находим $x_n: |x_n-a| < 1/n$, но $|f(x_n)-A| \ge \eps$.
Тогда $x_n \to a$, но $f(x_n) \not\to A$. Противоречие с Гейне.
\end{ticket}

% ================= БИЛЕТ 4.3 =================
\begin{ticket}{4.3 Критерий Коши (Функции)}
\textbf{Теорема:} $\exists \lim_{x\to a} f(x) \iff$
$\forall \eps>0 \exists \delta>0 \forall x', x'' \in \mathring{U}_\delta(a): |f(x') - f(x'')| < \eps$.
\tcblower
\textbf{Док-во (через Гейне):}
\textbf{($\Ra$)}: Неравенство $\triangle$. $|f(x')-f(x'')| \le |f(x')-A| + |f(x'')-A|$.
\textbf{($\Leftarrow$)}: Берем любую $x_n \to a$.
$x_n, x_m \in \mathring{U}_\delta \Ra |f(x_n)-f(x_m)| < \eps$.
$\{f(x_n)\}$ фундаментальна $\Ra$ сходится.
Все последовательности сходятся к одному пределу (смешанная посл-ть).
\end{ticket}

% ================= БИЛЕТ 4.4-4.5 =================
\begin{ticket}{4.4-4.5 Свойства пределов}
\textbf{1. Единственность.} (Через $\eps = |A-B|/2$ и неравенство треугольника $|A - B|$).
\textbf{2. Предел модуля:} $\lim f=A \Ra \lim |f|=|A|$.
\textbf{3. Неравенства:} $f(x) \le g(x) \Ra \lim f \le \lim g$.
(Док: От противного. $A > B \Ra$ неперес. окрестности).
\textbf{4. Зажатая функция:} $g \le f \le h, \lim g = \lim h = A \Ra \lim f = A$.
\textbf{5. Сохранение знака:} $\lim f = A \ne 0 \Ra \exists U_\delta: \sign f(x) = \sign A$.
(Док: $\eps = |A|/2$. $f(x)$ в $(A/2, 3A/2)$).
\textbf{6. Локальная огр-сть:} $\exists \lim \Ra \exists U_\delta$ где $f$ огр.
(Док: $\eps=1, |f| < |A|+1$).
\textbf{7. Арифметика:} Сумма, произв, частное.
(Док: через Гейне и св-ва последовательностей).
\end{ticket}

% ================= БИЛЕТ 4.6 =================
\begin{ticket}{4.6 Б.м. и Б.б. функции}
\textbf{Б.м.:} $\lim \alpha(x) = 0$. \textbf{Б.б.:} $\lim f(x) = \infty$.
\textbf{Связь:} $f(x) \to A \iff f(x) = A + \alpha(x)$.
\textbf{Св-ва б.м.:}
1. Сумма б.м. — б.м.
2. Б.м. $\times$ Ограниченная — б.м. ($|\alpha \cdot f| < \eps/M \cdot M$).
\textbf{Связь б.б./б.м.:}
$f \to \infty \Ra 1/f \to 0$.
$\alpha \to 0 (\alpha \ne 0) \Ra 1/\alpha \to \infty$.
\end{ticket}

% ================= БИЛЕТ 4.7-4.8 =================
\begin{ticket}{4.7-4.8 Замечательные пределы}
\textbf{1. Первый:} $\lim_{x\to 0} \frac{\sin x}{x} = 1$.
\textit{Док-во:} Геометрически ($x \in (0, \pi/2)$).
$S_{\triangle} < S_{sect} < S_{tan} \Ra \frac{1}{2}\sin x < \frac{1}{2}x < \frac{1}{2}\tg x$.
Делим на $\sin x$: $1 < \frac{x}{\sin x} < \frac{1}{\cos x}$. Зажимаем.
\textbf{2. Второй:} $\lim_{x\to \infty} (1+\frac{1}{x})^x = e$.
\textit{Док-во:} Сводим к посл-ти $n=[x]$.
$(1+\frac{1}{n+1})^n < (1+\frac{1}{x})^x < (1+\frac{1}{n})^{n+1}$.
Крайние $\to e$.
\end{ticket}

% ================= БИЛЕТ 4.9 =================
\begin{ticket}{4.9 Монотонные функции}
\textbf{Т:} Если $f(x) \uparrow$ на $(a, b)$, то существуют односторонние пределы:
$\lim_{x\to a+0} f(x) = \inf f(x)$, $\lim_{x\to b-0} f(x) = \sup f(x)$.
\textbf{Связь односторонних и двусторонних пределов:} если односторонние существуют и равно $\Leftarrow$ двусторонний равнен тому же.
\tcblower
\textbf{Док-во (для $x \to b-0$):}
Пусть $M = \sup f$.
1. $f(x) \le M$.
2. $\forall \eps \exists x_0: f(x_0) > M-\eps$.
Т.к. $f \uparrow$, для всех $x > x_0$: $f(x) \ge f(x_0) > M-\eps$.
Итог: $|f(x)-M| < \eps$.
\textbf{Следствие:} Разрывы только 1 рода (скачки).
\textbf{Доказательство связи:} Всё по определениям. 
\end{ticket}

% ================= БИЛЕТ 4.10-4.11 =================
\begin{ticket}{4.10-4.11 Непрерывность}
\textbf{В точке:} $\lim_{x\to a} f(x) = f(a)$.
\textbf{Коши:} $\forall \eps \exists \delta: |x-a|<\delta \Ra |f(x)-f(a)|<\eps$.
\textbf{Гейне:} $x_n \to a \Ra f(x_n) \to f(a)$.
\textbf{Арифметика:} Сумма, произв, частное непр. функций — непрерывны.
\textbf{Сложная ф-ция:} $t \to \tau, \phi(t) \to a, f(x) \to f(a)$.
Тогда $f(\phi(t)) \to f(a)$.
(Предел переходит через непрерывную функцию).
\end{ticket}

% ================= БИЛЕТ 4.12 =================
\begin{ticket}{4.12 Лемма Гейне-Бореля}
\textbf{Теорема:} Из любого открытого покрытия отрезка $[a, b]$ можно выделить \textbf{конечное} подпокрытие.
\tcblower
\textbf{Доказательство (От противного + Вложенные отрезки):}
1. Пусть нельзя выбрать конечное.
2. Делим $[a, b]$ пополам. Хотя бы одну половину нельзя покрыть конечным числом (иначе сумма конечных = конечное).
3. Повторяем деление. Получаем вложенные отрезки $\Delta_n \to 0$.
4. По Лемме Кантора есть общая точка $c \in \bigcap \Delta_n$.
5. Т.к. покрытие открытое, $c$ входит в некоторое множество $G_\alpha$ вместе с окрестностью $(c-\eps, c+\eps)$.
6. При большом $n$ отрезок $\Delta_n$ целиком попадет в эту окрестность.
7. Значит, $\Delta_n$ покрывается ОДНИМ множеством.
Противоречие (мы строили $\Delta_n$ так, что его нельзя покрыть конечным числом).
\end{ticket}

% ================= БИЛЕТ 4.13 =================
\begin{ticket}{4.13 Точки разрыва}
\textbf{1 рода:} Оба предела $\lim_{x\to a\pm 0} f(x)$ существуют и конечны.
- \textit{Устранимый:} $f(a-0) = f(a+0) \ne f(a)$. (Пример: $\frac{\sin x}{x}$ в 0).
- \textit{Скачок:} $f(a-0) \ne f(a+0)$. (Пример: $\sign x$).
\textbf{2 рода:} Хотя бы один предел не сущ. или $\infty$.
(Пример: $1/x$, $\sin(1/x)$).
\end{ticket}

% ================= БИЛЕТ 4.14 =================
\begin{ticket}{4.14 Теоремы Вейерштрасса (на отрезке)}
Пусть $f(x) \in C[a, b]$.
\textbf{1-я Теорема:} $f(x)$ ограничена на $[a, b]$.
\textit{Док-во:} От противного. Пусть неогр. $\exists x_n: |f(x_n)| > n$.
По Б-В выделим $x_{n_k} \to x_0$.
$f(x_{n_k}) \to f(x_0)$ (конечное). Но $|f| \to \infty$. Противоречие.
\textbf{2-я Теорема:} Достигает $\sup$ и $\inf$ ($\max$ и $\min$).
\textit{Док-во:} Пусть $M = \sup f$. Если не достигает, то $M - f(x) > 0$.
Рассм. $g(x) = \frac{1}{M-f(x)}$. Она непр., значит огр.
Но если $f(x)$ подходит к $M$, то $g(x) \to \infty$. Противоречие.
\end{ticket}

% ================= БИЛЕТ 4.15 =================
\begin{ticket}{4.15 Теорема о промеж. значении}
\textbf{Т (Больцано-Коши):} $f \in C[a, b]$ и $f(a)f(b) < 0 \Ra \exists c: f(c)=0$.
\textbf{Следствие:} Принимает все значения между $f(a)$ и $f(b)$.
\tcblower
\textbf{Док-во (Деление пополам):}
Строим вложенные отрезки $[a_n, b_n]$, на концах которых знаки разные.
$c = \lim a_n = \lim b_n$.
В силу непр., $f(c) = \lim f(a_n) \le 0$ и $f(c) = \lim f(b_n) \ge 0$.
Значит $f(c) = 0$.
\end{ticket}

% ================= БИЛЕТ 4.16 =================
\begin{ticket}{4.16 Теорема об обратной функции}
\textbf{Дано:} $f \in C[a, b]$, строго возрастает ($A=f(a), B=f(b)$).
\textbf{Утв:}
1. $E(f) = [A, B]$.
2. $\exists$ обратная $x = g(y)$, однозначная.
3. $g(y)$ строго возрастает.
4. $g(y)$ непрерывна на $[A, B]$.
\tcblower
\textbf{Док-во:}
1. $E(f)=[A,B]$: $f \uparrow \Ra A \le f \le B$. Непрерывна $\Ra$ принимает все промеж. значения.
2. Однозначность: $x_1 \ne x_2 \Ra f(x_1) \ne f(x_2)$ (строгая мон.).
3. Мон. $g$: Пусть $y_1 < y_2$. Если $x_1 \ge x_2 \Ra y_1 \ge y_2$ (противоречие).
4. Непр. $g$: Фиксируем $y_0, x_0$. Берем $\varepsilon$-окр $x_0$.
Концы $(x_0 \pm \varepsilon)$ переходят в $y_{\pm \varepsilon}$.
Интервал $(y_{-\varepsilon}, y_{+\varepsilon})$ является окрестностью $y_0$.
Внутри него выберем $\delta$-окр.
Её прообраз лежит внутри $(x_0-\varepsilon, x_0+\varepsilon)$.
\end{ticket}

% ================= БИЛЕТ 4.17 =================
\begin{ticket}{4.17 Критерий непр. монотонной}
\textbf{Теорема:} Пусть $f$ монотонна на $\Delta$.
$f$ непрерывна $\iff f(\Delta)$ — промежуток (без дырок).
\tcblower
\textbf{Док-во ($\Ra$):}
По теореме о промежуточном значении (Больцано-Коши).
Образ промежутка при непрерывном отображении — промежуток.
\textbf{Док-во ($\Leftarrow$):} (От противного).
Пусть образ — промежуток, но есть разрыв в $x_0$.
Т.к. $f$ монотонна, разрыв только 1 рода (скачок).
$\lim_{x\to x_0-0} f(x) < \lim_{x\to x_0+0} f(x)$.
Между пределами есть "дырка" (интервал значений), которые функция не принимает.
Значит $f(\Delta)$ не является сплошным промежутком. Противоречие.
\end{ticket}

% ================= БИЛЕТ 4.18 =================
\begin{ticket}{4.18 Критерий обратимости}
\textbf{Теорема:} Пусть $f \in C[a, b]$.
$f$ обратима $\iff f$ строго монотонна.
\tcblower
\textbf{($\Leftarrow$):} Строгая монотонность $\Ra$ разным $x$ разные $y$ $\Ra$ инъекция $\Ra$ обратима.
\textbf{($\Rightarrow$):} (От противного).
Пусть $f$ обратима, но не монотонна.
Тогда найдутся $x_1 < x_2 < x_3$ такие, что (например) $f(x_1) < f(x_2)$ и $f(x_2) > f(x_3)$ ("горка").
Возьмем $C$ чуть ниже вершины $f(x_2)$, но выше $\max(f(x_1), f(x_3))$.
По Т. о промеж. значении $f$ примет значение $C$ дважды:
один раз на $(x_1, x_2)$ и второй раз на $(x_2, x_3)$.
$f(c_1) = f(c_2) = C$.
Нарушена обратимость (инъективность). Противоречие.
\end{ticket}

% ================= БИЛЕТ 4.19 =================
\begin{ticket}{4.19 Равномерная непрерывность}
\textbf{Опр:} $\forall \eps \exists \delta \forall x', x'' (|x'-x''|<\delta \Ra |f(x')-f(x'')|<\eps)$.
($\delta$ не зависит от точки $x$).
\textbf{Т. Кантора:} $f \in C[a, b] \Ra f$ равномерно непрерывна.
\tcblower
\textbf{Док-во (Гейне-Бореля):}
Для каждой точки $x$ есть $\delta_x$ (локальная непр.).
Интервалы $(x-\delta_x/2, x+\delta_x/2)$ покрывают отрезок.
Выберем конечное покрытие.
Возьмем $\delta = \min (\delta_k/2)$.
Любые две близкие точки попадут в одну или соседние окрестности.
\end{ticket}

% ================= БИЛЕТ 4.20 =================
\begin{ticket}{4.20 O-символика}
$f = o(g) \iff f/g \to 0$ (б.м. высшего порядка).
$f = O(g) \iff |f/g| \le C$ (ограничена).
$f \sim g \iff f/g \to 1$ (эквивалентны).
\textbf{Главная часть:} $f(x) = A(x-a)^k + o(\dots)$.
\textbf{Таблица ($x \to 0$):}
$\sin x \sim x$, $\tg x \sim x$, $\ln(1+x) \sim x$, $e^x-1 \sim x$.
$1-\cos x \sim x^2/2$, $(1+x)^\alpha - 1 \sim \alpha x$.
\end{ticket}

% ================= БИЛЕТ 4.21-4.22 =================
\begin{ticket}{4.21-4.22 Рост функций}
\textbf{Иерархия роста ($x \to \infty$):}
$\ln x \ll x^\alpha \ll a^x \ll n! \ll n^n$ ($\alpha > 0, a > 1$).
\textbf{Пределы:}
1. $\lim \frac{x^\alpha}{a^x} = 0$. (Док: $n/b^n \to 0$ по Биному Ньютона).
2. $\lim \frac{\ln x}{x^\alpha} = 0$. (Док: замена $x=e^t \to t/e^{\alpha t} \to 0$).
3. $\lim_{x\to 0} x^\alpha \ln x = 0$. (Док: $x=1/t \to -\ln t / t^\alpha \to 0$).
\end{ticket}



% ================= БИЛЕТ 5.1 =================
\begin{ticket}{5.1 Производная и Дифф-сть}
\textbf{Опр:} $f'(x_0) = \lim_{\Delta x \to 0} \frac{f(x_0+\Delta x) - f(x_0)}{\Delta x}$.
\textbf{Односторонние:} $f'_{\pm}(x_0) = \lim_{\Delta x \to \pm 0} \dots$.
\textbf{Дифф-сть:} $\Delta f = A \Delta x + o(\Delta x)$.
\textbf{Связь:} Дифф-ма $\iff \exists f'(x_0)$. При этом $A = f'(x_0)$.
\tcblower
\textbf{Док-во связи:}
($\Ra$) $\frac{\Delta f}{\Delta x} = A + \frac{o(\Delta x)}{\Delta x} \to A$.
($\Leftarrow$) $\lim \frac{\Delta f}{\Delta x} = f' \Ra \frac{\Delta f}{\Delta x} = f' + \alpha \Ra \Delta f = f'\Delta x + \alpha \Delta x$.
\end{ticket}

% ================= БИЛЕТ 5.2 =================
\begin{ticket}{5.2 Дифференциал. Геометрия}
\textbf{Дифференциал:} $df = f'(x_0) dx$ (главная линейная часть приращения).
\textbf{Геом. смысл $f'$:} Тангенс угла наклона касательной ($k = \tg \alpha$).
\textbf{Геом. смысл $df$:} Приращение ординаты касательной.
\textbf{Связь:} Дифф-ма $\Ra$ Непрерывна ($\Delta f \to 0$).
\tcblower
\textbf{Уравнение касательной:} $y - f(x_0) = f'(x_0)(x - x_0)$.
\end{ticket}

% ================= БИЛЕТ 5.3 =================
\begin{ticket}{5.3 Арифметика производных}
$(u \pm v)' = u' \pm v'$.
$(uv)' = u'v + uv'$.
$(u/v)' = \frac{u'v - uv'}{v^2}$.
\tcblower
\textbf{Док-во (произведения):}
$\Delta(uv) = u(x+\Delta x)v(x+\Delta x) - u(x)v(x)$.
Добавим и вычтем $u(x)v(x+\Delta x)$.
$\dots = u(x) \Delta v + v(x+\Delta x) \Delta u$.
Делим на $\Delta x$, переходим к пределу.
\end{ticket}

% ================= БИЛЕТ 5.4 =================
\begin{ticket}{5.4 Таблица производных}
1. $(C)' = 0$.
2. $(x^\alpha)' = \alpha x^{\alpha-1}$.
3. $(\sin x)' = \cos x$, $(\cos x)' = -\sin x$.
4. $(e^x)' = e^x$, $(a^x)' = a^x \ln a$.
5. $(\ln x)' = 1/x$, $(\log_a x)' = \frac{1}{x \ln a}$.
\tcblower
\textbf{Вывод $\ln x$:}
$\frac{\ln(x+\Delta x) - \ln x}{\Delta x} = \frac{1}{\Delta x} \ln(1+\frac{\Delta x}{x}) = \ln(1+\frac{\Delta x}{x})^{\frac{1}{\Delta x}}$.
Замена $t = \Delta x/x$. Степень $\to 1/x \cdot 1/t$.
Предел $(1+t)^{1/t} = e$. Итог $\frac{1}{x} \ln e = \frac{1}{x}$.
\end{ticket}

% ================= БИЛЕТ 5.5 =================
\begin{ticket}{5.5 Сложная и Обратная}
\textbf{Сложная:} $(f(\varphi(x)))' = f'_y \cdot \varphi'_x$.
\textbf{Обратная:} $(f^{-1})'_y = \frac{1}{f'_x}$.
\textbf{Арки:}
$(\arcsin x)' = \frac{1}{\sqrt{1-x^2}}$.
$(\arccos x)' = -\frac{1}{\sqrt{1-x^2}}$.
$(\arctg x)' = \frac{1}{1+x^2}$.
\tcblower
\textbf{Док-во (Обратная):}
$\Delta x / \Delta y = 1 / (\Delta y / \Delta x)$. Предел дроби = $1/f'$.
\textbf{Вывод arcsin:} $x = \sin y \Ra x'_y = \cos y$.
$y'_x = 1/\cos y = 1/\sqrt{1-\sin^2 y} = 1/\sqrt{1-x^2}$.
\end{ticket}

% ================= БИЛЕТ 5.6 =================
\begin{ticket}{5.6 Лейбниц (n-я производная)}
\textbf{Формула:} $(uv)^{(n)} = \sum_{k=0}^n C_n^k u^{(k)} v^{(n-k)}$.
Похожа на бином Ньютона.
\tcblower
\textbf{Док-во:} Индукция.
Шаг $n \to n+1$: дифференцируем сумму.
$(u^{(k)} v^{(n-k)})' = u^{(k+1)} v^{(n-k)} + u^{(k)} v^{(n-k+1)}$.
Группируем, используем св-во $C_n^k + C_n^{k-1} = C_{n+1}^k$.
\end{ticket}

% ================= БИЛЕТ 5.7 =================
\begin{ticket}{5.7 Ферма и Ролль}
\textbf{Ферма:} $x_0$ — экстремум, $\exists f' \Ra f'(x_0) = 0$.
\textit{Док-во:} $\Delta f = \Delta x (f' + \alpha)$. Если $f' > 0$, знак меняется $\Ra$ не экстремум.
\textbf{Ролль:} $f \in C[a,b], \exists f' \in (a,b), f(a)=f(b) \Ra \exists \xi: f'(\xi)=0$.
\textit{Док-во:} По Вейерштрассу есть $\max$ и $\min$.
Если равны $\Ra$ const. Если разные $\Ra$ экстр. внутри $\Ra$ Ферма.
\end{ticket}

% ================= БИЛЕТ 5.8 =================
\begin{ticket}{5.8 Лагранж и Коши}
\textbf{Лагранж:} $\frac{f(b)-f(a)}{b-a} = f'(\xi)$.
\textbf{Коши:} $\frac{f(b)-f(a)}{g(b)-g(a)} = \frac{f'(\xi)}{g'(\xi)}$.
\tcblower
\textbf{Док-во (Лагранж):} Всп. ф-ция $F(x) = f(x) - \frac{f(b)-f(a)}{b-a}(x-a)$.
$F(a)=F(b) \Ra$ Ролль $\Ra F'(\xi)=0$.
\textbf{Док-во (Коши):} $F(x) = f(x) - \frac{f(b)-f(a)}{g(b)-g(a)} g(x)$.
\end{ticket}

% ================= БИЛЕТ 5.9 =================
\begin{ticket}{5.9 Постоянство и Монотонность}
\textbf{Const:} $f'(x) \equiv 0 \iff f = \text{const}$.
\textbf{Монотонность:}
$f'(x) \ge 0 \iff f \uparrow$.
$f'(x) \le 0 \iff f \downarrow$.
$f'(x) > 0 \Ra f \uparrow\uparrow$ (строго).
\tcblower
\textbf{Док-во:} Всё через Лагранжа: $f(x_2)-f(x_1) = f'(\xi)(x_2-x_1)$.
Знак разности зависит от знака производной.
\end{ticket}

% ================= БИЛЕТ 5.10-5.11 =================
\begin{ticket}{5.10-5.11 Тейлор (Лагранж / Пеано)}
$f(x) = \sum_{k=0}^n \frac{f^{(k)}(a)}{k!} (x-a)^k + R_n$.
\textbf{Лагранж (Глобал):} $R_n = \frac{f^{(n+1)}(\xi)}{(n+1)!}(x-a)^{n+1}$.
(Док-во: многократный Коши или интегрирование по частям).
\textbf{Пеано (Локал):} $R_n = o((x-a)^n)$ при $x \to a$.
(Док-во: Правило Лопиталя $n-1$ раз).
\end{ticket}

% ================= БИЛЕТ 5.12 =================
\begin{ticket}{5.12 Тейлор для $e^x, \sin, \cos$}
$e^x = 1 + x + x^2/2! + \dots + x^n/n! + R_n$.
$\sin x = x - x^3/3! + x^5/5! - \dots$.
$\cos x = 1 - x^2/2! + x^4/4! - \dots$.
\textbf{Оценка остатка:}
$|R_n| \le \frac{M |x|^{n+1}}{(n+1)!}$.
Т.к. $M$ огр ($e^\xi$ или 1), а факториал растет быстрее степени $\Ra R_n \to 0$ для всех $x$.
\end{ticket}





\begin{comment}

% ================= БИЛЕТ 3.9 (Задачи) =================
\begin{ticket}{3.9 (Доп) Задачи про число e}
\textbf{1. Связь с рядом:} $e = \sum \frac{1}{k!}$.
\textit{Док-во:} Бином Ньютона для $(1+1/n)^n$.
Каждое слагаемое $\to 1/k!$. Сумма ограничена $e$.
\textbf{2. Монотонность $(1+1/n)^{n+p}$:}
Исследуем $\ln f(x)$. Производная $\approx \frac{1}{x^2}(\frac{1}{2}-p)$.
$p < 1/2 \Ra \uparrow$. $p > 1/2 \Ra \downarrow$.
$p=1/2 \Ra$ убывает быстрее всего ($\approx -1/6x^3$).
\textbf{3. Точность $10^{-3}$ (Ряд):}
Остаток $R_n < \frac{1}{n \cdot n!}$.
$n=6$: $1/4320 < 0.001$. Ответ: $n=6$.
\textbf{4. Точность $10^{-3}$ (Последовательность):}
Ошибка $\approx \frac{e}{2n}$.
$\frac{e}{2n} < 0.001 \Ra n > 500e \approx 1360$.
\end{ticket}

\end{comment}


\end{multicols*}
\end{document}
